\section{The Mesh}

The whole of space is the mesh, one growing honeycomb made of identical cells joined face to face, filling a curved space with no gaps. Its shape has a name in the standard notation for regular honeycombs, the four-dimensional hyperbolic $\{3,4,3,4\}$, whose cells are 24-cells. Three facts carry the weight. It is discrete, a countable weave of identical cells. It is hyperbolic, curved so that room opens up faster than in flat space. And it grows.

The mesh is built from the bottom up, and the order is forced. A tone is one value at one site. Twenty-four sites close into one dock, the 24-cell of the previous step. Each of the dock's twenty-four sites points outward through a face, so a single dock cannot stand alone, it demands a neighbor across every face, twenty-four identical docks pressed around it. Those neighbors demand their own neighbors in turn, and the wake lays the new docks down at the open edge beat after beat. The mesh is the running total, one dock copied face to face without end. Tone, then dock, then mesh, each one forced by the one before it.

Why a honeycomb at all. Because the base must be discrete, and a regular honeycomb is the simplest discrete space there is, every cell alike, every joint alike, no special places and no edges to fall off.

Why hyperbolic rather than flat. A flat lattice picks out preferred directions and breaks the even-handedness a relativistic world must have. Hyperbolic space has no preferred direction, and it has room. In flat space a region grows only as fast as its surface, but in hyperbolic space the interior swells, tree-like, with far more room a few steps out than any flat lattice allows. That extra room is what later makes the boundary of a region behave like a screen, and lets a bounded region hold a high-dimensional world inside itself.

The shape is committed to rather than assumed, on a selection argument. Among the regular hyperbolic honeycombs, which form a short finite list because they run out by the fifth dimension, demand three things at once: that it be crystallographic, that its cell carry spinors, and that its flat boundary be three-dimensional. The decisive test is the spinors. The twenty-four directions of $\{3,4,3,4\}$ are the $D_4$ roots, which carry spinors, while the three-dimensional sibling $\{5,3,4\}$ has twelve directions that provably carry none, so spin would have to be imposed by hand there. $\{3,4,3,4\}$ is the unique honeycomb that hands over spin for free and has a three-dimensional boundary, so it is selected. The lower siblings appear in this paper only to build intuition and as controls, never as the subject.

\begin{figure}[ht]
\centering
\includegraphics[width=0.5\textwidth]{vibe-mesh-73}
\caption{A $\{7,3\}$ tessellation of the hyperbolic plane on the Poincar\'e disk, seven-sided cells fanning out from every point with room to branch. This flat two-dimensional picture is only an aid to the imagination, a stand-in we can actually draw, since the substrate the theory runs on is the four-dimensional hyperbolic honeycomb $\{3,4,3,4\}$, which cannot be pictured directly. The three colors mark the three values of the tone, minus one, zero, and plus one.}
\label{fig:tessellation-73}
\end{figure}

\begin{figure}[ht]
\centering
\includegraphics[width=\textwidth]{vibe-mesh-534}
\caption{A view inside the $\{5,3,4\}$ honeycomb, the three-dimensional hyperbolic packing of dodecahedral cells. Like the last figure it is a stand-in for the imagination, the closest cousin we can render in three dimensions, since the committed substrate is the four-dimensional $\{3,4,3,4\}$, whose 24-cell honeycomb cannot be pictured directly. This is the dimension-down sibling used as a control, the one whose twelve directions carry no spinor.}
\label{fig:tessellation-534}
\end{figure}
